This thesis investigates the security of human-chosen 6-digit PINs under low-entropy attacks. 
Although a 6-digit PIN has a theoretical space of one million combinations, user behavior 
introduces strong biases that significantly reduce practical security.

This work models human-chosen PIN distributions and evaluates multiple attack strategies, 
including random guessing, frequency-based guessing, and personal-information-assisted attacks. 
The study focuses on limited-attempt scenarios such as top-1, top-3, top-5, and top-10 guesses.

Experimental results demonstrate that user bias and personal information leakage greatly increase 
attack success rates compared to the uniform random model. The thesis highlights the importance 
of guessability-based metrics and provides recommendations to improve PIN security.